\section{Logarithmic derivatives and explicit series}

For an input equipped with pole separation define
$\psi(z)=\Gamma'(z)/\Gamma(z)$; this division is licensed by
$|\Gamma(z)|\geq1/M_R$.  Differentiating the finite factorization (C7)
and its explicitly convergent power series gives
\[
 \boxed{\displaystyle
 \psi(z)=-\gamma-\frac1z+
       \sum_{k=1}^{\infty}\frac{z}{k(k+z)}.}
 \tag{C27}
\]
For $N\geq2R$ the tail after $N$ is at most $2R/N$.
Indeed $|k+z|\geq k/2$, and therefore the $k$th term is at most
$2R/k^2$.  This also explains why the separated form is useful: each
finite head term has a certified denominator, and the tail has a
uniform bound independent of the fine pole separation.

An equivalent partial fraction series is
\[
 \psi(z)=-\gamma+
 \sum_{k=0}^{\infty}\left(\frac1{k+1}-\frac1{k+z}\right).
 \tag{C28}
\]
The difference of corresponding finite sums is an explicitly vanishing
endpoint term.  For every integer $m\geq1$,
\[
 \boxed{\displaystyle
 \psi^{(m)}(z)=(-1)^{m+1}m!
               \sum_{k=0}^{\infty}(k+z)^{-m-1}.}
 \tag{C29}
\]
If $N\geq2R$, the tail over $k>N$ has absolute value at most
\[
 \frac{(m-1)!2^{m+1}}{N^m}.
 \tag{C30}
\]
These differentiations are not justified merely by convergence of
(C27).  On a slightly larger pole-separated disk, differentiating
$L(z/k)$ $r$ times gives, for $r\geq2$,
\[
 \frac{d^r}{dz^r}L(z/k)
   =\frac{(-1)^{r-1}(r-1)!}{(k+z)^r},
\]
with tail bounded by
$(r-1)!2^r/[(r-1)N^{r-1}]$.
For $r=1$ its tail is at most $2R/N$.
Finite Taylor remainders, bounded by one further derivative on that
larger disk, prove the derivative assertions with explicit moduli.
All lengths, factorials, and tail budgets remain primitive recursive
uniformly in the requested derivative order.

Differentiating the functional equation or reindexing finite sums gives
\[
 \psi(z+1)=\psi(z)+z^{-1},\qquad
 \psi^{(m)}(z+1)=\psi^{(m)}(z)+(-1)^m m!z^{-m-1}.
 \tag{C31}
\]
In particular, for $n\geq0$ and $m\geq1$,
\[
 \begin{aligned}
 \psi(n+1)&=H_n-\gamma,\\
 \psi^{(m)}(1)&=(-1)^{m+1}m!\zeta(m+1),\\
 \psi^{(m)}(n+1)&=(-1)^{m+1}m!
       \bigl(\zeta(m+1)-H_n^{(m+1)}\bigr).
 \end{aligned}
 \tag{C32}
\]
Separation of even and odd terms in (C29), a finite reindexing followed
by (C30), yields
\[
 \psi^{(m)}(1/2)=(-1)^{m+1}m!(2^{m+1}-1)\zeta(m+1).
 \tag{C33}
\]
Similarly the $N$th partial sum in (C28) at $z=1/2$ is
$-\gamma+2H_N-2H_{2N}$.  By (C2),
$H_{2N}-H_N-\log2=\gamma_{2N}-\gamma_N$ has absolute value at most
$1/N$.  Hence
\[
 \psi(1/2)=-\gamma-2\log2.
 \tag{C34}
\]
For positive real $x$ the series in (C29) has positive terms, giving
the quantitative sign statement
\[
 (-1)^{m+1}\psi^{(m)}(x)\geq m!x^{-m-1}>0.
 \tag{C35}
\]
In particular $\psi$ is strictly increasing and $\log\Gamma$ is
strictly convex on the positive real axis.  On a rational interval
$[a,b]$ the lower bound $1/b^2$ for $\psi'$ supplies a modulus for
strict convexity; it is stronger than a purely qualitative assertion.

Higher gamma derivatives can be computed without new analytic arguments:
\[
 \Gamma^{(n)}(z)=\Gamma(z)Y_n(z),\quad Y_0=1,\qquad
 Y_{n+1}=\sum_{j=0}^n\binom nj Y_{n-j}\psi^{(j)}.
 \tag{C36}
\]
This is the ordinary finite Leibniz rule applied to
$\Gamma'=\Gamma\psi$, and is a primitive recursive polynomial
recursion for every derivative order.

\section{Taylor and Laurent expansions with remainders}

For $|w|\leq r<1$, (C11), (C20), and finite geometric bounds prove
\[
 \Gamma(1+w)=\exp\!\left(
 -\gamma w+\sum_{m=2}^{\infty}
                    \frac{(-1)^m\zeta(m)}m w^m\right).
 \tag{C37}
\]
The cancellation of the factor $w$ is cancellation of power-series
coefficients (or of the corresponding finite polynomial factors),
so this argument also applies at $w=0$ without deciding whether an
arbitrary named complex number is zero.
The exponent is the local branch of $\log\Gamma(1+w)$ whose value at
$w=0$ is zero.  The wording is essential: no global principal logarithm
identity is asserted.  Since $1\leq\zeta(m)\leq2$ for $m\geq2$,
the tail after degree $d\geq1$ of this exponent is at most
\[
 \frac{2r^{d+1}}{(d+1)(1-r)}.
 \tag{C38}
\]
The same argument, or (C27), gives
\[
 \psi(1+w)=-\gamma+
       \sum_{m=1}^{\infty}(-1)^{m+1}\zeta(m+1)w^m,
 \tag{C39}
\]
whose tail after degree $d$ is at most $2r^{d+1}/(1-r)$.
To turn these into explicit primitive recursive degree bounds, choose
an integer $q\geq2$ with $r\leq1-1/q$.  The finite binomial inequality
gives $r^q\leq1/2$, so a degree of the form $q(p+c)$ reduces the
geometric factor below $2^{-p-c}$.  The remaining prefactors are
absorbed by an explicitly chosen integer $c$.

Writing $\Gamma(1+w)=\sum_{n\geq0}g_nw^n$, finite coefficient
multiplication gives
\[
 g_0=1,\qquad
 ng_n=-\gamma g_{n-1}+
         \sum_{m=2}^n(-1)^m\zeta(m)g_{n-m}.
 \tag{C40}
\]
For example
\[
 g_1=-\gamma,\quad
 g_2=\frac{\gamma^2+\zeta(2)}2,\quad
 g_3=-\frac{\gamma^3+3\gamma\zeta(2)+2\zeta(3)}6.
\]
Thus $\Gamma'(1)=-\gamma$ and
$\Gamma''(1)=\gamma^2+\zeta(2)$.
If a direct Taylor-polynomial error is desired on $|w|\leq r<1$,
choose a rational $r<s<1$.  The coefficient absolute sum at $s$ is at
most
\[
 A_s=\exp\left(s+\frac{s^2}{1-s}\right),
\]
because $2/m\leq1$ for $m\geq2$.  Hence
\[
 \left|\Gamma(1+w)-\sum_{n=0}^d g_nw^n\right|
 \leq A_s(r/s)^{d+1}.
 \tag{C41}
\]
An integer power of $3$ can replace $A_s$ to make the majorant rational.

Writing $P(z)=\sum_{n\geq1}c_nz^n$, one obtains
\[
 c_1=1,\quad c_2=\gamma,\qquad
 (n-1)c_n=\gamma c_{n-1}+
    \sum_{m=2}^{n-1}(-1)^{m+1}\zeta(m)c_{n-m}\quad(n\geq3).
 \tag{C42}
\]
For example
\[
 c_3=\frac{\gamma^2-\zeta(2)}2,\qquad
 c_4=\frac{\gamma^3-3\gamma\zeta(2)+2\zeta(3)}6.
\]
Although its derivation starts in $|z|<1$, this reciprocal Maclaurin
series converges on every disk by (C19), with the same coefficients.

There is also a concrete Laurent remainder at every pole.  For
$0<|h|\leq1/4$, combine (C24) and (C37) to write
\[
 \Gamma(-n+h)=\frac{(-1)^n}{n!h}\exp\bigl(Ah+V(h)\bigr),
 \quad A=H_n-\gamma,
\]
where
\[
 V(h)=\sum_{m=2}^{\infty}
 \frac{(-1)^m\zeta(m)+H_n^{(m)}}m h^m,
 \qquad |V(h)|\leq3|h|^2.
\]
Since $|A|\leq n+1$, the exponential remainder gives
\[
 \left|\Gamma(-n+h)-\frac{(-1)^n}{n!}
       \left(\frac1h+H_n-\gamma\right)\right|
 \leq\frac{C_n}{n!}|h|,
 \quad C_n=3+3^{n+2}(n+2)^2.
 \tag{C43}
\]
The generous integer constant avoids any issue of computing an
exponential bound.  At $n=0$ this includes
$\Gamma(h)=h^{-1}-\gamma+O(h)$ with the displayed effective meaning.

\subsection{Taylor series at every positive integer}
The same cancellation gives a wider family of branch-free exponential
formulas.  For $n\geq1$ and $|w|<n$,
\[
 \Gamma(n+w)=(n-1)!\exp\!\left((H_{n-1}-\gamma)w+
   \sum_{m=2}^{\infty}\frac{(-1)^m}{m}
          \bigl(\zeta(m)-H_{n-1}^{(m)}\bigr)w^m\right).
 \tag{C44}
\]
This follows either by cancelling the first $n-1$ factors of (C11),
or by multiplying (C37) by $\prod_{j=1}^{n-1}(j+w)$ and comparing
finite coefficients.  The former proof immediately supplies the full
disk $|w|<n$, without an analytic continuation theorem.  Since
\[
 0<\zeta(m)-H_{n-1}^{(m)}
       \leq n^{-m}+\frac{n^{1-m}}{m-1}\leq(n+1)n^{-m},
\]
the exponent tail after degree $d\geq1$, on $|w|\leq r<n$, is at most
\[
 \frac{(n+1)(r/n)^{d+1}}{(d+1)(1-r/n)}.
 \tag{C45}
\]
For $n=2$ this is the branch-free version of the commonly quoted
log-gamma expansion involving $\zeta(m)-1$.  If it is rewritten as a
formula for $\log\Gamma(1+w)$, the additional logarithm of $1+w$
requires a compatible local branch; (C44) avoids that issue entirely.

\section{Two further finite-product consequences}

Let $x,y$ be real and supply pole separation for $x$.  Taking absolute
values in the finite Euler quotient gives the exact finite identity
\[
 \left|\frac{\Gamma_N(x)}{\Gamma_N(x+iy)}\right|^2
       =\prod_{k=0}^N\left(1+\frac{y^2}{(x+k)^2}\right).
\]
Here $|N^{iy}|=1$, an identity of the elementary exponential.
The left side has a specified limit by (C15).  Alternatively, the
logarithm of the right-hand tail is bounded by $4y^2/N$ as soon as
$N\geq2|x|$, since $\log(1+t)\leq t$ for $t\geq0$.
Thus
\[
 \boxed{\displaystyle
 \left|\frac{\Gamma(x)}{\Gamma(x+iy)}\right|^2
       =\prod_{k=0}^{\infty}\left(1+\frac{y^2}{(x+k)^2}\right).}
 \tag{C46}
\]
In particular $|\Gamma(x+iy)|\leq|\Gamma(x)|$, and if a positive
lower bound for $|y|$ is supplied, the inequality is strict with a
quantitative margin already furnished by the $k=0$ factor.  No
positivity assumption on $x$ is needed, only its separation from the
nonpositive integers.

More generally, let finite lists $a_1,\ldots,a_m$ and
$b_1,\ldots,b_m$ satisfy $\sum_i a_i=\sum_i b_i$.  Supply pole
separation for the $b_i$; the $a_i$ need no such condition.  Finite
cancellation proves
\[
 \prod_{k=0}^{N}\prod_{i=1}^m\frac{a_i+k}{b_i+k}
       =\prod_{i=1}^m\frac{Q_N(a_i)}{Q_N(b_i)}.
\]
Hence
\[
 \boxed{\displaystyle
 \prod_{k=0}^{\infty}\prod_{i=1}^m\frac{a_i+k}{b_i+k}
       =\prod_{i=1}^m P(a_i)\Gamma(b_i).}
 \tag{C47}
\]
If the $a_i$ are also separated from the poles, the right side is
$\prod_i\Gamma(b_i)/\Gamma(a_i)$.  Allowing $P(a_i)$ instead makes
the zero cases meaningful as well.  Suppose all the parameters have
absolute value at most $R$ and all $b_i$ have separation $\delta$.
The proof of (C14) also bounds $|Q_N|$ by $M_R$.  The difference of
one quotient in the last finite display and its limit is therefore
at most $6M_R^2B_R/(\lambda_{R,\delta}^2N)$.
The finite telescoping product identity yields the explicit total error
\[
 \frac{6mM_R^2B_R}{\lambda_{R,\delta}^2N}
       \left(\frac{M_R}{\lambda_{R,\delta}}\right)^{m-1}
 \quad(m\geq1).
 \tag{C48}
\]
The empty-list case is the exact equality $1=1$.

\paragraph{Reference alignment and formalization boundary.}
The identities in this development agree with NIST DLMF
\S5.2 (definitions and poles), \S5.7 (series), \S5.8 (products), and
\S5.15 (polygamma functions).  The proof above supplies its own finite
estimates.  Every infinite statement is read as its family of rational
error inequalities.  A Lean theorem stated directly over Mathlib's
$\mathbb C$ is a useful presentation theorem but is not, by itself,
an iPRA proof.  The arithmetic layer must verify the finite products,
tail bounds, coefficient recursions, and precision transformers, with
quantifier-free induction on their displayed finite recursions.
